解 \(Ax=0\) 时，两个解相加还是解。解 \(y''+y=0\) 时，两个解相加还是解。次数不超过 2 的多项式相加，次数还是不超过 2。把两张同尺寸的图逐像素平均，得到的还是一张同尺寸的图。四个场合的对象毫不相干，验证却是同一行代数。

一段推理在这么多互不相干的地方原样重复，说明它不属于其中任何一个。本章把它单独拎出来：向量空间就是把``可以线性组合''这条性质从具体对象中剥离出来的产物。剥离之后，凡满足这条性质的集合都共用同一套关于方向、维数与分解的几何直觉——哪怕它的元素既没有箭头，也没有天然的坐标。

为此要重新分清两样东西。对象是对象，坐标是描述。选一组基相当于挑一把尺子，同一个向量在两把尺子下读数完全不同，被读的东西却纹丝未动；换基改的是账本，不是被记的操作。这个区分一旦立住，后面``为什么要找特征向量''、``为什么要正交化''、``为什么 SVD 值得算''就都不再是技巧问题，而是同一个问题的不同答案：挑哪把尺子，能让结构自己显形。

\begin{center}
\begin{tikzpicture}[x=1cm,y=1cm,font=\footnotesize,>=stealth]
  \draw[black!50,fill=black!4] (0,2.4) rectangle (2.6,3.3);
  \node[align=center] at (1.3,2.85) {\(\R^n\) 中的向量\\齐次方程的解};
  \draw[black!50,fill=black!4] (3.2,2.4) rectangle (5.8,3.3);
  \node[align=center] at (4.5,2.85) {多项式、函数\\微分方程的解};
  \draw[black!50,fill=black!4] (6.4,2.4) rectangle (9.0,3.3);
  \node[align=center] at (7.7,2.85) {图像、信号\\嵌入向量};
  \draw[->,black!60] (1.3,2.35) -- (2.6,1.85);
  \draw[->,black!60] (4.5,2.35) -- (4.5,1.85);
  \draw[->,black!60] (7.7,2.35) -- (6.4,1.85);
  \draw[black!55,fill=blue!10] (1.6,0.95) rectangle (7.4,1.8);
  \node at (4.5,1.38) {共用的唯一性质：可以线性组合};
  \draw[->,black!60] (3.2,0.9) -- (2.4,0.35);
  \draw[->,black!60] (5.8,0.9) -- (6.6,0.35);
  \draw[black!55] (0.2,-0.65) rectangle (4.2,0.3);
  \node[align=center] at (2.2,-0.17) {基与坐标：一次描述选择\\维数是不动的那个数};
  \draw[black!55] (4.8,-0.65) rectangle (8.8,0.3);
  \node[align=center] at (6.8,-0.17) {四个子空间：解是否存在\\是否唯一、冗余在哪};
  \draw[->,black!60] (2.2,-0.7) -- (3.4,-1.3);
  \draw[->,black!60] (6.8,-0.7) -- (5.6,-1.3);
  \draw[black!55,fill=blue!10] (1.6,-2.3) rectangle (7.4,-1.35);
  \node[align=center] at (4.5,-1.83) {换基：描述改变，对象不变\\同一变换的矩阵彼此相似};
\end{tikzpicture}

\emph{抽象一次，四处兑现。中间那条性质是全章唯一用到的假设；上面是它覆盖的对象，下面是它换来的两件工具，最后合成一条判断：矩阵里的数字有一半属于操作，另一半只属于你选的坐标。}
\end{center}

四篇沿这条线推进。\hyperref[sec:03-Linear-Algebra/02-Vector-Spaces/01]{``为什么需要向量空间''}交代八条公理各自挡住什么失效，并给出子空间的单条判据：几何上就是``必须穿过原点''，平移一下就只剩仿射集。\hyperref[sec:03-Linear-Algebra/02-Vector-Spaces/02]{``基、坐标与维数：用最少的方向完整描述空间''}把``够用''与``不多余''这两个相反的要求压在一组向量上，说明坐标依赖基而维数不依赖，后者正是实践中所说的有效自由度。

\hyperref[sec:03-Linear-Algebra/02-Vector-Spaces/03]{``矩阵的四个基本子空间：输入、输出与丢失的信息''}是本章分量最重的一篇，也是全书反复回看的一张图：输入空间被行空间与零空间垂直劈开，输出空间被列空间与左零空间垂直劈开，中间只有 \(r\) 维能穿过去。方程何时有解、解是否唯一、参数为什么不可辨识、最小范数解为什么落在行空间里，都由这张图一次回答。\hyperref[sec:03-Linear-Algebra/02-Vector-Spaces/04]{``线性变换与换基：对象不变，坐标可以改变''}最后把操作与它的账本分开，给出 \(R^{-1}AS\) 与相似变换，并说明哪些量在换基下岿然不动。

阅读顺序上，前两篇是后两篇的语言基础，第一次读按顺序走最省力。若眼下的困难是``模型系数不稳定''、``方程组无解或多解''这类具体症状，可以直接翻到第三篇，再往回补概念。动手前应当熟悉 \hyperref[ch:016]{``线性方程组与矩阵''}里的消元与主元——本章几乎每个结论都要靠主元个数 \(r\) 来落地。

读完之后能带走的判断是这样几条：拿到一个矩阵，不必算就能说出它丢掉多少自由度、什么样的右端项才可能有解；看到两个数字迥异的矩阵，会先问它们是不是同一个变换换了基；再听到``换个特征表示效果就变好了''，知道变的是抽取效率而不是信息量。本章不涉及长度与角度，因此还答不出``哪一组基最好''——那需要内积，是下一单元的事。
